% This file is included by uet_thesis.tex.
% Replace blue placeholders and remove example text before submission.

\chapter{Project Overview}
\label{chap:overview}
% =============================================================================

\section{Executive Summary}
\begin{guidance}
Summarize the problem, intended users, proposed approach, principal deliverables,
and expected value. Write this section last and keep it understandable without
the rest of the document.
\end{guidance}

\placeholder{Provide a concise executive summary, normally 150--250 words.}

\section{Background and Motivation}
\begin{guidance}
Explain the context, current practice, affected users or organizations, and why
the problem is worth addressing. Support factual claims with credible references.
Do not repeat the complete literature review.
\end{guidance}

\placeholder{Describe the project context and motivation.}

\section{Problem Statement}
\begin{guidance}
State the existing condition, the specific gap or difficulty, its consequences,
and the need for the proposed work. Describe the problem rather than presenting a
chosen technology as the problem.
\end{guidance}

\begin{examplebox}
Students currently arrange advising appointments through unstructured messages.
This makes availability difficult to track, causes scheduling conflicts, and
provides no reliable appointment record. A structured mechanism is therefore
needed to coordinate and record advising appointments.
\end{examplebox}

\placeholder{State the project problem in one focused paragraph.}

\section{Objectives}
\begin{guidance}
Define the specific outcomes the project intends to achieve. Objectives should be
clear, feasible, and measurable. Describe intended results rather than routine
activities such as ``conduct a literature review'' or technologies such as
``use Flutter.'' Each objective must be supported by scope items, requirements,
design elements, and verification evidence through Chapter~\ref{chap:planning}.
\end{guidance}

\begin{table}[H]
\centering
\caption{Project objectives}
\label{tab:objectives}
\begin{tabularx}{\textwidth}{@{}p{0.14\textwidth}X@{}}
\toprule
\textbf{ID} & \textbf{Objective} \\
\midrule
OBJ-01 & Enable registered students to request and track advising appointments through a single system. \\
OBJ-02 & \placeholder{Specific and measurable project outcome} \\
\bottomrule
\end{tabularx}
\end{table}

\section{Proposed Approach}
\begin{guidance}
Provide a high-level description of how the project will address the stated
problem. Development and game projects should summarize the proposed solution;
research projects should summarize the research approach; hybrid projects should
state both. Detailed architecture and methodology belong in Chapter~\ref{chap:design}.
\end{guidance}

\begin{examplebox}
The project will develop a web-based appointment portal with separate student
and adviser roles. A hybrid project could additionally evaluate whether the
portal reduces scheduling time through a controlled user study.
\end{examplebox}

\placeholder{Describe the proposed solution and/or research approach.}

\section{Scope}
\begin{guidance}
Define the approved boundaries of the project. Each scope item should express one
capability, user group, platform, data source, integration, or deliverable and
should map to at least one objective. Detailed behaviour belongs in the
requirements chapter.
\end{guidance}

\subsection{In Scope}

\begin{table}[H]
\centering
\caption{In-scope items}
\label{tab:in-scope}
\begin{tabularx}{\textwidth}{@{}p{0.14\textwidth}X p{0.20\textwidth}@{}}
\toprule
\textbf{Scope ID} & \textbf{Included item} & \textbf{Objective(s)} \\
\midrule
SCP-01 & Appointment requests by registered students. & OBJ-01 \\
SCP-02 & Appointment approval and rescheduling by advisers. & OBJ-01 \\
SCP-03 & \placeholder{One project capability or boundary} & \placeholder{OBJ-XX} \\
\bottomrule
\end{tabularx}
\end{table}

\subsection{Out of Scope \applicability{If applicable}}
\begin{guidance}
State intentional exclusions and their reasons. If there are no meaningful
exclusions, state this explicitly rather than inventing them.
\end{guidance}

\begin{table}[H]
\centering
\caption{Out-of-scope items}
\label{tab:out-of-scope}
\begin{tabularx}{\textwidth}{@{}p{0.14\textwidth}X X@{}}
\toprule
\textbf{ID} & \textbf{Excluded item} & \textbf{Reason} \\
\midrule
EXC-01 & Integration with the university's central timetable system. & The external interface is unavailable during the FYP period. \\
EXC-02 & \placeholder{Excluded item} & \placeholder{Clear boundary reason} \\
\bottomrule
\end{tabularx}
\end{table}

\section{Users and Stakeholders}
\begin{guidance}
Identify direct users and other parties who influence, maintain, fund, regulate,
or benefit from the project. State their main needs and influence on the design.
\end{guidance}

\begin{table}[H]
\centering
\caption{Users and stakeholders}
\label{tab:stakeholders}
\begin{tabularx}{\textwidth}{@{}p{0.22\textwidth}p{0.18\textwidth}X X@{}}
\toprule
\textbf{Stakeholder} & \textbf{Role} & \textbf{Need or interest} & \textbf{Design influence} \\
\midrule
Registered student & Primary user & Request and track appointments & Simple mobile-responsive interface \\
\placeholder{Stakeholder} & \placeholder{Role} & \placeholder{Need} & \placeholder{Design implication} \\
\bottomrule
\end{tabularx}
\end{table}

% =============================================================================
